\documentclass[11pt]{article}

\usepackage[margin=1in]{geometry}
\usepackage{amsmath, amssymb}
\usepackage{hyperref}
\usepackage{enumitem}
\usepackage{xcolor}
\usepackage{listings}

\lstdefinestyle{julia}{
  basicstyle=\ttfamily\small,
  breaklines=true,
  frame=single,
  columns=fullflexible,
  keywordstyle=\color{blue!60!black},
  commentstyle=\color{green!40!black},
  stringstyle=\color{orange!60!black}
}
\lstset{style=julia}

\title{Notes on the Architecture, Moves, Energies, and Metropolis Rule in \texttt{ParticlesMC}}
\author{}
\date{}

\begin{document}
\maketitle

\section{Purpose of the Code}

\texttt{ParticlesMC} is a Julia package for Monte Carlo simulations of particle
and molecular systems. The package is built as a domain-specific layer on top of
\texttt{Arianna}. In this architecture, \texttt{Arianna} provides the generic
Monte Carlo simulation machinery, while \texttt{ParticlesMC} defines:

\begin{itemize}[nosep]
  \item what a particle or molecular system is;
  \item how pair and bonded energies are computed;
  \item which Monte Carlo moves are available;
  \item how proposed moves are applied and reverted;
  \item how configurations and trajectories are loaded and stored;
  \item how a user-facing TOML file is converted into a simulation.
\end{itemize}

\section{High-Level Architecture}

The main package entry point is:

\begin{center}
\texttt{src/ParticlesMC.jl}
\end{center}

This file defines the module, imports \texttt{Arianna}, includes the other source
files, exports the public API, and defines the command-line interface
\texttt{particlesmc}.

\subsection{Main Source Files}

\begin{description}[leftmargin=3.2cm, style=nextline]
  \item[\texttt{ParticlesMC.jl}]
  Top-level module, exports, generic helpers, and the \texttt{particlesmc}
  command-line interface.

  \item[\texttt{atoms.jl}]
  Defines the \texttt{Atoms} system type for non-bonded particle systems and its
  energy calculation.

  \item[\texttt{molecules.jl}]
  Defines the \texttt{Molecules} system type for systems with molecular identity
  and bonds. It contains both bonded and non-bonded energy calculations.

  \item[\texttt{orientation.jl}]
  Defines replaceable molecular-orientation conventions, periodic molecular
  unwrapping, and mass-weighted center-of-mass calculations.

  \item[\texttt{models.jl}]
  Defines interaction models and potentials, for example Lennard-Jones,
  smoothed Lennard-Jones, soft spheres, and generalized Kremer-Grest-type
  interactions.

  \item[\texttt{neighbours.jl}]
  Defines neighbor-list data structures used to reduce the cost of local energy
  calculations.

  \item[\texttt{moves.jl}]
  Defines Monte Carlo actions such as particle displacements, species swaps, and
  molecule flips. It also defines proposal policies and methods required by
  \texttt{Arianna}.

  \item[\texttt{utils.jl}]
  Defines the target Boltzmann weight, periodic boundary helpers, species-list
  helpers, and the energy callback.

  \item[\texttt{IO/IO.jl}]
  Defines loading and storing of configurations. Format-specific logic is split
  into \texttt{IO/xyz.jl}, \texttt{IO/exyz.jl}, and \texttt{IO/lammps.jl}.

  \item[\texttt{rotation.jl}]
  Defines a custom \texttt{AriannaAlgorithm} for molecular rotation observables.
\end{description}

\subsection{Simulation Flow}

The command-line path is:

\begin{enumerate}[nosep]
  \item read a TOML parameter file;
  \item read initial configurations through \texttt{load\_chains};
  \item build either \texttt{Atoms} or \texttt{Molecules} systems;
  \item build the interaction model matrix;
  \item construct a pool of moves;
  \item construct an \texttt{Arianna} algorithm list;
  \item create \texttt{Simulation(chains, algorithm\_list, steps)};
  \item call \texttt{run!(simulation)}.
\end{enumerate}

Conceptually:

\[
\text{TOML input}
\rightarrow
\text{systems}
\rightarrow
\text{moves + algorithms}
\rightarrow
\text{Arianna simulation}
\rightarrow
\text{outputs}.
\]

\section{System Types}

\subsection{Particles}

The abstract base type is:

\begin{lstlisting}[language=Julia]
abstract type Particles <: AriannaSystem end
\end{lstlisting}

Both \texttt{Atoms} and \texttt{Molecules} are subtypes of \texttt{Particles}.
This lets generic methods dispatch on \texttt{Particles}, while specialized
methods can dispatch on \texttt{Atoms} or \texttt{Molecules}.

\subsection{Atoms}

\texttt{Atoms} represents a non-bonded particle system. Important fields include:

\begin{itemize}[nosep]
  \item \texttt{position}: particle positions;
  \item \texttt{species}: particle species labels;
  \item \texttt{density}, \texttt{temperature};
  \item \texttt{energy}: total system energy, stored as a one-element vector;
  \item \texttt{model\_matrix}: pair interaction model for every species pair;
  \item \texttt{box}: periodic simulation box;
  \item \texttt{neighbour\_list}: object used to find local neighbors;
  \item \texttt{species\_list}: helper used for species-swap moves.
\end{itemize}

\subsection{Molecules}

\texttt{Molecules} extends the particle-system idea with molecular structure:

\begin{itemize}[nosep]
  \item \texttt{molecule}: molecule id for each particle;
  \item \texttt{mass}: mass of each particle;
  \item \texttt{start\_mol}, \texttt{length\_mol}: start and length of each molecule;
  \item \texttt{Nmol}: number of molecules;
  \item \texttt{bonds}: bonded neighbors for each particle;
  \item \texttt{Phi}: rotation-vector observable storage.
\end{itemize}

Mass is stored once per particle. The public \texttt{System} constructor accepts
the optional keyword \texttt{masses}. If it is omitted, every particle receives
unit mass, preserving the behavior of configurations created before mass support
was introduced. Supplied masses must be finite, strictly positive, and have the
same length as the position array. For command-line simulations, masses may be
specified in the \texttt{[system]} TOML table:

\begin{lstlisting}
[system]
masses = [1.0, 2.0, 1.0]
\end{lstlisting}

Mass does not modify the configurational pair Hamiltonian or the current Monte
Carlo acceptance rule. It is presently used to define molecular centers of mass
and orientations.

\section{Molecular Orientation}

Molecular orientation is implemented separately in
\texttt{src/orientation.jl}. This separation allows the geometrical convention
to change without modifying Monte Carlo moves or the future external-field
energy term. All orientation functions return a unit vector and account for
periodic boundaries by first placing the atoms in mutually consistent images.

For positions \(\mathbf r_i\) and masses \(m_i\), the center of mass is

\begin{equation}
\mathbf r_{\mathrm{COM}}
=
\frac{\sum_i m_i\mathbf r_i}{\sum_i m_i}.
\end{equation}

It is available through:

\begin{lstlisting}[language=Julia]
center = molecule_center_of_mass(system, molecule_id)
\end{lstlisting}

\subsection{Center-to-Atom Convention}

\texttt{CenterToAtomOrientation(a)} defines the orientation from the center of
mass to local atom \(a\):

\begin{equation}
\mathbf n
=
\frac{\mathbf r_a-\mathbf r_{\mathrm{COM}}}
     {\lVert\mathbf r_a-\mathbf r_{\mathrm{COM}}\rVert}.
\end{equation}

For example:

\begin{lstlisting}[language=Julia]
definition = CenterToAtomOrientation(1)
n = orientation(definition, system, molecule_id)
\end{lstlisting}

The atom index is local to the selected molecule rather than a global particle
index.

\subsection{Triangular Plane-Normal Convention}

\texttt{PlaneNormalOrientation(a,b,c)} defines the right-handed normal to the
surface determined by three ordered local atoms:

\begin{equation}
\mathbf n
=
\frac{(\mathbf r_b-\mathbf r_a)\times
      (\mathbf r_c-\mathbf r_a)}
     {\lVert(\mathbf r_b-\mathbf r_a)\times
      (\mathbf r_c-\mathbf r_a)\rVert}.
\end{equation}

The default is \texttt{PlaneNormalOrientation(1,2,3)}. Exchanging two indices
reverses \(\mathbf n\), so atom ordering fixes the sign needed by a polar field
coupling. The plane normal itself is independent of the atomic masses. Collinear
or coincident defining coordinates are rejected because they do not define a
unique normal.

\subsection{Aligning External Field}

The aligning field is represented by \texttt{AligningField(h, definition)} and
adds the one-body molecular energy

\begin{equation}
H_{\mathrm{field}}
=
-\sum_{m=1}^{N_{\mathrm{mol}}}\mathbf h\cdot\mathbf n_m.
\end{equation}

It is intentionally separate from pair-interaction models. During system
initialization, the pair contribution is still divided by two, while each field
term is added exactly once. A system without a field uses
\texttt{NoExternalField()}, preserving previous behavior.

An aligning field can be specified in TOML as:

\begin{lstlisting}
[system.external_field]
type = "align"
h = [0.0, 0.0, 0.7]
orientation = "plane_normal"
atoms = [1, 2, 3]
\end{lstlisting}

For the center-to-atom convention, use
\texttt{orientation = "center\_to\_atom"} and \texttt{atom = 1}.

For an ordinary particle displacement, only the orientation of the displaced
particle's molecule is recomputed. The additional trial energy is

\begin{equation}
\Delta U_h
=
-\mathbf h\cdot
\left(\mathbf n_m^{\mathrm{new}}-\mathbf n_m^{\mathrm{old}}\right).
\end{equation}

The implementation dispatches field scoring on the action type. Only
\texttt{Displacement} is field-aware at present. A future geometric flip move
can provide its own action-specific method because its affected orientation and
energy accounting need not be identical to a one-particle displacement. The
current \texttt{MoleculeFlip}, which swaps species labels, is unchanged.

\section{Interaction Models and Hamiltonians}

The interaction models are defined in \texttt{src/models.jl}. The base types are:

\begin{lstlisting}[language=Julia]
abstract type Model end
abstract type DiscreteModel <: Model end
abstract type ContinuousModel <: Model end
\end{lstlisting}

For a pair of particles \(i,j\), the code obtains the appropriate interaction by
species:

\[
M_{ij} = \texttt{model\_matrix}_{s_i,s_j}.
\]

Then it evaluates either:

\[
U_{ij}^{\mathrm{nb}} = V_{M_{ij}}(r_{ij}^2),
\]

or, for bonded molecular pairs,

\[
U_{ij}^{\mathrm{b}} = V_{M_{ij}}^{\mathrm{bond}}(r_{ij}^2).
\]

Here \(r_{ij}\) is computed using the nearest-image convention for periodic
boundary conditions.

\subsection{Hamiltonian for Particle Systems}

For an \texttt{Atoms} system, the Hamiltonian is purely pairwise and non-bonded:

\begin{equation}
H_{\mathrm{atoms}}(\mathbf{r}, \mathbf{s})
=
\sum_{1 \le i < j \le N}
V_{s_i s_j}(r_{ij}^2)
\mathbf{1}\!\left(r_{ij}^2 \le r_{\mathrm{cut},s_i s_j}^2\right).
\end{equation}

Equivalently, the code computes a local energy:

\begin{equation}
E_i
=
\sum_{j \in \mathcal{N}(i)}
V_{s_i s_j}(r_{ij}^2),
\end{equation}

where \(\mathcal{N}(i)\) is supplied by the chosen neighbor list. The total
energy is initialized as:

\begin{equation}
H_{\mathrm{atoms}} = \frac{1}{2}\sum_{i=1}^N E_i.
\end{equation}

The factor \(1/2\) appears because a pair \(i,j\) is counted once in \(E_i\) and
once in \(E_j\).

\subsection{Hamiltonian for Molecular Systems}

For a \texttt{Molecules} system, the Hamiltonian has bonded and non-bonded parts:

\begin{equation}
H_{\mathrm{mol}}
=
H_{\mathrm{bond}}
+
H_{\mathrm{nonbond}}.
\end{equation}

The bonded part is:

\begin{equation}
H_{\mathrm{bond}}
=
\sum_{(i,j)\in \mathcal{B}}
V_{s_i s_j}^{\mathrm{bond}}(r_{ij}^2),
\end{equation}

where \(\mathcal{B}\) is the set of bonds.

The non-bonded part is:

\begin{equation}
H_{\mathrm{nonbond}}
=
\sum_{\substack{1 \le i < j \le N \\ (i,j)\notin \mathcal{B}}}
V_{s_i s_j}(r_{ij}^2)
\mathbf{1}\!\left(r_{ij}^2 \le r_{\mathrm{cut},s_i s_j}^2\right).
\end{equation}

Locally, for one particle \(i\), the code computes:

\begin{equation}
E_i
=
\sum_{j\in \mathcal{B}(i)}
V_{s_i s_j}^{\mathrm{bond}}(r_{ij}^2)
+
\sum_{j\in \mathcal{N}(i),\, j\notin \mathcal{B}(i)}
V_{s_i s_j}(r_{ij}^2).
\end{equation}

Then the initial total energy is again:

\begin{equation}
H_{\mathrm{mol}} = \frac{1}{2}\sum_{i=1}^N E_i.
\end{equation}

\section{Where Energy is Computed in the Code}

The generic entry point is in \texttt{src/ParticlesMC.jl}:

\begin{lstlisting}[language=Julia]
function compute_energy_particle(system::Particles, i::Int)
    return compute_energy_particle(system, i, system.neighbour_list)
end
\end{lstlisting}

\subsection{Atoms Energy}

For \texttt{Atoms}, the local particle energy is computed in
\texttt{src/atoms.jl}:

\begin{lstlisting}[language=Julia]
function compute_energy_particle(system::Atoms, i, neighbour_list::NeighbourList)
    energy_i = zero(typeof(system.density))
    position_i = get_position(system, i)
    for j in neighbour_list(system, i)
        energy_i += compute_energy_ij(system, i, j, position_i)
    end
    return energy_i
end
\end{lstlisting}

The pair energy is:

\begin{lstlisting}[language=Julia]
function compute_energy_ij(system::Atoms, i, j, position_i)
    i == j && return zero(typeof(system.density))
    model_ij = get_model(system, i, j)
    position_j = get_position(system, j)
    r2 = nearest_image_distance_squared(position_i, position_j, get_box(system))
    r2 > cutoff2(model_ij) && return zero(typeof(system.density))
    return potential(r2, model_ij)
end
\end{lstlisting}

\subsection{Molecules Energy}

For \texttt{Molecules}, local energy is computed in \texttt{src/molecules.jl}:

\begin{lstlisting}[language=Julia]
function compute_energy_particle(system::Molecules, i, neighbour_list::NeighbourList)
    position_i = system.position[i]
    bonds_i = system.bonds[i]

    energy_i = compute_energy_bonded_i(system, i, position_i, bonds_i)
    for j in neighbour_list(system, i)
        energy_i += check_nonbonded_compute_energy_ij(system, i, j, position_i, bonds_i)
    end
    return energy_i
end
\end{lstlisting}

Bonded energy calls:

\begin{lstlisting}[language=Julia]
bond_potential(r2, model_ij)
\end{lstlisting}

Non-bonded energy calls:

\begin{lstlisting}[language=Julia]
potential(r2, model_ij)
\end{lstlisting}

\section{Moves and the \texttt{Arianna} Interface}

\texttt{src/moves.jl} defines the available Monte Carlo actions and how they
interact with the \texttt{Arianna} engine.

\subsection{Important Distinction}

\texttt{moves.jl} does not implement the full Metropolis algorithm. Instead, it
implements methods that \texttt{Arianna}'s \texttt{Metropolis} algorithm calls.

That division is:

\begin{center}
\begin{tabular}{ll}
\textbf{\texttt{Arianna}} & Metropolis loop, accept/reject logic, scheduling \\
\textbf{\texttt{ParticlesMC}} & moves, energy differences, proposal densities, reverts
\end{tabular}
\end{center}

\subsection{The Central Hook}

The main method connecting moves to \texttt{Arianna} is:

\begin{lstlisting}[language=Julia]
function Arianna.perform_action!(system::Particles, action::Action)
    e1, e2 = perform_action!(system, action)
    if isinf(e1) || isinf(e2)
        action.delta_e = zero(typeof(system.energy[1]))
    else
        action.delta_e = e2 - e1
        system.energy[1] += action.delta_e
    end
    return e1, e2
end
\end{lstlisting}

In the source code the field is written as \texttt{delta e} using the Julia
unicode symbol \texttt{delta e}; in these notes we write it as \(\Delta E\).

The purpose is:

\begin{enumerate}[nosep]
  \item compute the energy before the proposed move, \(E_1\);
  \item apply the move;
  \item compute the energy after the move, \(E_2\);
  \item store \(\Delta E = E_2 - E_1\);
  \item temporarily update the system energy.
\end{enumerate}

If \texttt{Arianna} rejects the move, it calls the corresponding
\texttt{Arianna.revert\_action!} method.

\subsection{Displacement Move}

\texttt{Displacement} moves one particle:

\[
\mathbf{r}_i' = \mathbf{r}_i + \boldsymbol{\delta}.
\]

The proposal policy \texttt{SimpleGaussian} samples:

\[
i \sim \mathrm{Uniform}\{1,\dots,N\},
\qquad
\boldsymbol{\delta} \sim \mathcal{N}(0,\sigma^2 I).
\]

Because the Gaussian displacement is symmetric,

\[
q(x\to x') = q(x'\to x),
\]

so the proposal-density terms cancel in ordinary Metropolis sampling.

The code computes only the local energy of particle \(i\) before and after the
move:

\[
\Delta E = E_i(\mathbf{r}_i') - E_i(\mathbf{r}_i).
\]

\subsection{Discrete Species Swap}

\texttt{DiscreteSwap} swaps the species labels of two particles:

\[
s_i' = s_j,
\qquad
s_j' = s_i.
\]

The policy \texttt{DoubleUniform} selects one particle from each of two species
lists:

\[
i \sim \mathrm{Uniform}(\{k:s_k=a\}),
\qquad
j \sim \mathrm{Uniform}(\{k:s_k=b\}).
\]

The log proposal density in the code is:

\[
\log q(x\to x') =
-\log(N_a N_b),
\]

where \(N_a\) and \(N_b\) are the numbers of particles of the two selected
species.

The energy change is computed from the local energies of particles \(i\) and
\(j\):

\[
\Delta E =
\left(E_i' + E_j'\right)
-
\left(E_i + E_j\right).
\]

\subsection{Energy-Biased Swap}

\texttt{EnergyBias} is a non-uniform proposal policy. It samples particles with
weights depending on their local energies:

\[
w_i^{(a)} \propto \exp(\theta_1 E_i),
\qquad
w_j^{(b)} \propto \exp(\theta_2 E_j).
\]

The proposal density is therefore not generally symmetric. This is why the code
implements:

\begin{lstlisting}[language=Julia]
Arianna.log_proposal_density(action, ::EnergyBias, parameters, system)
\end{lstlisting}

For non-symmetric proposals, the Metropolis-Hastings correction must include
the proposal-ratio term.

\subsection{Molecule Flip}

\texttt{MoleculeFlip} swaps two sites within the same molecule:

\[
s_i' = s_j,
\qquad
s_j' = s_i,
\qquad
i,j \in \text{same molecule}.
\]

It is similar to \texttt{DiscreteSwap}, but constrained by molecular membership.

\section{Where the Metropolis Rule Is}

The Metropolis algorithm itself is not defined in \texttt{src/moves.jl}. It is
provided by \texttt{Arianna}. In this package, \texttt{Metropolis} appears when
building the algorithm list:

\begin{lstlisting}[language=Julia]
(algorithm=Metropolis, pool=pool, seed=seed,
 parallel=parallel, sweepstep=length(chains[1]))
\end{lstlisting}

The Boltzmann target density needed by Metropolis is defined in
\texttt{src/utils.jl}:

\begin{lstlisting}[language=Julia]
function Arianna.unnormalised_log_target_density(e, system::Particles)
    return -e / system.temperature
end
\end{lstlisting}

and:

\begin{lstlisting}[language=Julia]
function Arianna.delta_log_target_density(e1, e2, system::Particles)
    return -(e2 - e1) ./ system.temperature
end
\end{lstlisting}

Thus, for a symmetric proposal, the acceptance probability is:

\begin{equation}
P_{\mathrm{acc}}(x\to x')
=
\min\left[1,\exp\left(-\frac{\Delta E}{T}\right)\right].
\end{equation}

For a non-symmetric proposal, the Metropolis-Hastings form is:

\begin{equation}
P_{\mathrm{acc}}(x\to x')
=
\min\left[
1,
\exp\left(
-\frac{\Delta E}{T}
+
\log q(x'\to x)
-
\log q(x\to x')
\right)
\right].
\end{equation}

The proposal-density terms are supplied by \texttt{moves.jl} through methods
like:

\begin{lstlisting}[language=Julia]
Arianna.log_proposal_density(action, policy, parameters, system)
\end{lstlisting}

\section{Summary}

\begin{itemize}
  \item The full simulation loop and Metropolis accept/reject rule are provided
  by \texttt{Arianna}.
  \item \texttt{ParticlesMC} defines systems, energies, moves, proposal policies,
  and file IO.
  \item Energy is computed locally using neighbor lists.
  \item For \texttt{Atoms}, the Hamiltonian is a sum of non-bonded pair
  potentials.
  \item For \texttt{Molecules}, the Hamiltonian contains both bonded and
  non-bonded contributions.
  \item \texttt{moves.jl} tells \texttt{Arianna} how to sample, apply, score, and
  revert moves.
  \item The Boltzmann target used by Metropolis is \(\log \pi(x) = -H(x)/T\).
\end{itemize}

\end{document}
